\section*{ТЕРМИНЫ И ОПРЕДЕЛЕНИЯ}

В настоящей работе используются следующие термины с соответствующими
определениями.

\begin{description}\setlength{\itemsep}{0.2em}
\item[Автономная система (AS)] -- совокупность сетевых префиксов,
управляемых одним административным субъектом и анонсируемых наружу
по единой маршрутной политике; идентифицируется числом ASN
(Autonomous System Number).
\item[BGP (Border Gateway Protocol)] -- протокол обмена маршрутной
информацией между автономными системами, описанный в RFC~4271
\cite{rfc4271}.
\item[BMP (BGP Monitoring Protocol)] -- протокол наблюдения за
состоянием BGP-сеансов и содержимым таблиц маршрутизации
\cite{rfc7854}.
\item[CIDR (Classless Inter-Domain Routing)] -- модель бесклассовой
адресации, в которой принадлежность адреса префиксу описывается парой
«адрес/длина маски» \cite{rfc4632}.
\item[Длиннейшее совпадение префикса (Longest Prefix Match, LPM)] --
задача поиска по IP-адресу префикса наибольшей длины, которому этот
адрес принадлежит.
\item[Полный снимок маршрутов (snapshot)] -- единомоментный набор
действующих BGP-маршрутов, формирующих полное начальное состояние
таблицы.
\item[Поток обновлений (stream)] -- последовательность отдельных
BGP-сообщений типа update/withdraw, изменяющих состояние таблицы
после загрузки начального снимка.
\item[Adaptive Radix Tree (ART)] -- адаптивное префиксное дерево с
переменным размером узлов, описанное в работе Лейса и др.
\cite{leis2013art}.
\item[\texttt{ranger}] -- прежняя внутренняя реализация хранения
префиксных диапазонов, использовавшаяся для построения таблицы
«IP-префикс~$\rightarrow$~ASN» в BIFIT COLLECTOR. В работе
\texttt{ranger} рассматривается как базовый вариант для сравнения с
новым шардированным ART-хранилищем.
\item[GoBGP] -- открытая реализация BGP-демона на языке Go,
предоставляющая внешнее gRPC-API для управления и наблюдения.
\item[gRPC] -- система удалённого вызова процедур поверх HTTP/2,
используемая для взаимодействия микросервиса с GoBGP и с клиентами.
\item[BIFIT MITIGATOR] -- продукт компании, обеспечивающий защиту от
распределённых атак отказа в обслуживании.
\item[BIFIT COLLECTOR] -- продукт для сбора и анализа сетевого
трафика крупного оператора связи.
\item[SLA (Service Level Agreement)] -- соглашение об уровне
обслуживания, фиксирующее количественные требования к качеству работы
сервиса.
\end{description}

%==============================================================
